% preamble_summary.tex: packages, theorem environments and general macros of the abridged version.
% Copied from the preamble of the long version (main.tex); the draft-note machinery is omitted.
% GENUS workshop style: page geometry, fonts, hyperref, fancyhdr and the claim-status tags (genus-collab/meta)
\usepackage{genus}
% text width of the original note (170 mm), so that its displayed formulas still fit the line
\geometry{margin=2cm}
\usepackage[utf8]{inputenc}
\usepackage[T1]{fontenc}

%-------------------------------------------------------------------------------
% Packages (modelled on the preamble of [BS] and of the companion paper [LP2]; only what is needed)
\usepackage[dvipsnames]{xcolor}
\usepackage{amssymb}
\usepackage[mathscr]{eucal}
\usepackage{amsthm}
\usepackage{amsmath}
\usepackage{mathtools}
\definecolor{highlight}{HTML}{3465a4}
% hyperref is loaded by genus.sty, with the GENUS link colours
\hypersetup{hypertexnames=false, hyperfootnotes=false, linktocpage=true, bookmarksdepth=2}
\usepackage{cleveref}
\usepackage{url}
\usepackage{array}

%-------------------------------------------------------------------------------
% Typefaces. [BS] uses libertine with newtxmath; libertine needs fontaxes, which is not assumed
% here, so mathpazo (Palatino) is the substitute, as in the companion paper.
% (mathpazo and courier dropped: genus.sty sets Latin Modern)
\usepackage{microtype}
\linespread{1.06}
\setlength{\parskip}{4pt}
\setlength{\parindent}{0pt}
\setcounter{secnumdepth}{3}
% enumitem is not assumed: first-level enumerate is labelled (i), (ii), ... ; for named parts use
% description or \item[(F1)] with an explicit label.
\renewcommand{\labelenumi}{(\roman{enumi})}
\renewcommand{\theenumi}{(\roman{enumi})}
\numberwithin{equation}{section}

%-------------------------------------------------------------------------------
% Theorem environments (names as in [BS])
\theoremstyle{plain}
\newtheorem{thm}{Theorem}[section]
\newtheorem{lem}[thm]{Lemma}
\newtheorem{prop}[thm]{Proposition}
\newtheorem{cor}[thm]{Corollary}
\theoremstyle{definition}
\newtheorem{defn}[thm]{Definition}
\newtheorem{rmk}[thm]{Remark}
\newtheorem{notation}[thm]{Notation}
% genus.sty defines an example environment on its own counter; redefine it on the counter of thm
\let\example\undefined\let\endexample\undefined\expandafter\let\csname c@example\endcsname\undefined
\newtheorem{example}[thm]{Example}
\newtheorem{conv}[thm]{Convention}

\crefname{equation}{Eq.}{Eqs.}
\crefname{thm}{Theorem}{Theorems}
\crefname{lem}{Lemma}{Lemmas}
\crefname{prop}{Proposition}{Propositions}
\crefname{cor}{Corollary}{Corollaries}
\crefname{defn}{Definition}{Definitions}
\crefname{rmk}{Remark}{Remarks}
\crefname{notation}{Notation}{Notations}
\crefname{example}{Example}{Examples}
\crefname{conv}{Convention}{Conventions}
\crefname{section}{Section}{Sections}
\crefname{subsection}{Section}{Sections}
\crefname{appendix}{Appendix}{Appendices}

%-------------------------------------------------------------------------------
% Letters
\newcommand{\bbA}{\mathbb{A}}
\newcommand{\bbC}{\mathbb{C}}
\newcommand{\bbD}{\mathbb{D}}
\newcommand{\bbE}{\mathbb{E}}
\newcommand{\bbF}{\mathbb{F}}
\newcommand{\bbK}{\mathbb{K}}
\newcommand{\bbL}{\mathbb{L}}
\newcommand{\bbP}{\mathbb{P}}
\newcommand{\bbQ}{\mathbb{Q}}
\newcommand{\bbR}{\mathbb{R}}
\newcommand{\bbT}{\mathbb{T}}
\newcommand{\bbZ}{\mathbb{Z}}
\newcommand{\cA}{\mathcal{A}}
\newcommand{\cC}{\mathcal{C}}
\newcommand{\cE}{\mathcal{E}}
\newcommand{\cF}{\mathcal{F}}
\newcommand{\cG}{\mathcal{G}}
\newcommand{\cL}{\mathcal{L}}
\newcommand{\cN}{\mathcal{N}}
\newcommand{\cO}{\mathcal{O}}
\newcommand{\cP}{\mathcal{P}}
\newcommand{\cQ}{\mathcal{Q}}
\newcommand{\cW}{\mathcal{W}}
\newcommand{\cZ}{\mathcal{Z}}
\newcommand{\mgp}[1]{\mathsf{#1}}

%-------------------------------------------------------------------------------
% General
\newcommand{\dd}{\mathrm{d}}
\newcommand{\ri}{\mathrm{i}}
\newcommand{\re}{\mathrm{e}}
\newcommand{\de}{\partial}
% \llbracket: from newtxmath if available, otherwise a substitute (stmaryrd is not assumed)
\providecommand{\llbracket}{[\![}
\providecommand{\rrbracket}{]\!]}
\newcommand{\fps}[1]{\llbracket #1 \rrbracket}
\newcommand{\lser}[1]{(\!(#1)\!)}
\newcommand{\Frac}{\mathrm{Frac}}
\DeclareMathOperator{\Tr}{Tr}
\DeclareMathOperator{\End}{End}
\DeclareMathOperator{\Aut}{Aut}
\DeclareMathOperator{\Spec}{Spec}
\DeclareMathOperator*{\Res}{Res}
\DeclareMathOperator{\Disc}{Disc}
\DeclareMathOperator{\Li}{Li}
\DeclareMathOperator{\val}{val}
\DeclareMathOperator{\ev}{ev}
\newcommand{\tensor}{\otimes}
\newcommand{\coef}[1]{[#1]}                    % coefficient extraction, e.g. \coef{t^k}

%-------------------------------------------------------------------------------
% Geometry, tori, weights (weights as in [BS]: \alpha_i, \epsilon_i)
\newcommand{\Tmax}{\mgp{T}}                    % the five-torus
\newcommand{\TCY}{\mgp{T}_{\mathrm{CY}}}        % the Calabi--Yau four-torus E = \epsilon_1 + \epsilon_2
\newcommand{\TF}{\mgp{T}_{\mathrm{F}}}          % fibrewise torus \alpha_0 = 0, \alpha_i = \epsilon_i
\newcommand{\Tthree}{\mgp{T}_{3}}               % \alpha_0 = 0 inside \TCY
\newcommand{\wA}{\mathsf{A}}                    % \alpha_0 + \alpha_2 = E - \alpha_1
\newcommand{\wB}{\mathsf{B}}                    % \alpha_0 + \alpha_1 = E - \alpha_2
\newcommand{\uq}{\mathsf{u}}                    % (E - 2\alpha_1)(E - 2\alpha_2)
\newcommand{\Mbar}{\overline{M}}
\newcommand{\vir}{\mathrm{virt}}
\newcommand{\GW}{\mathsf{GW}}
\newcommand{\Hdg}{\bbE}                        % Hodge bundle
\newcommand{\Lamv}{\Lambda^{\!\vee}}           % \Lamv(u) = \sum_i (-1)^i \lambda_i u^{g-i}
\newcommand{\Hg}{\mathsf{Hg}}                  % \Hg_{\bw} = \prod_a \Lamv(w_a)/w_a^g (normalised Hodge class)
\newcommand{\Nb}{\cN_{\mathrm{cf}}}             % \cO(-1)_{\alpha_1} + \cO(-1)_{\alpha_2} + \cO_{-\epsilon_1} + \cO_{-\epsilon_2}
\newcommand{\bw}{\mathbf{w}}                   % (E,-\epsilon_1,-\epsilon_2)
\newcommand{\nw}{\mathsf{nw}}                  % (\alpha_1,\wA,-\epsilon_1,-\epsilon_2)
\newcommand{\wts}{\mathrm{wt}}                 % \wts_\pm = the five tangent weights at p_\pm
\newcommand{\eul}{\mathrm{e}}                  % \eul(\wts_\pm): product of the tangent weights
\newcommand{\sF}{\mathscr{F}}                  % \sF_g(\cQ) = \sum_{d\ge 0} \GW_{g,d}\cQ^d
\newcommand{\sG}{\mathscr{G}}                  % \sG_g = \sF_g - \sF_g(0)

%-------------------------------------------------------------------------------
% References to the long version of this paper: \Lref{Theorem~4.1} prints [L, Theorem 4.1].
\newcommand{\Lref}[1]{[L,~#1]}

